\documentclass[a4paper,12pt]{report}
\usepackage[table]{xcolor}
\usepackage{tabularx}
\usepackage{tabularray}
\usepackage{array}
%\usepackage[french]{babel}
\usepackage[utf8]{inputenc}
\usepackage[T1]{fontenc}
\usepackage{lmodern}
\usepackage[a4paper,left=2cm,right=2cm,top=2.5cm,bottom=2.5cm]{geometry}
\usepackage{graphicx}
\usepackage{colortbl}
\usepackage{hhline}
\usepackage[hidelinks]{hyperref}
\usepackage{setspace}
\usepackage{titlesec}
\usepackage{tikz}

% --- Palette technique ---
\definecolor{hdr}{HTML}{1A2E3C}
\definecolor{rowA}{HTML}{EFF4F7}
\definecolor{rowB}{HTML}{FFFFFF}
\definecolor{accent}{HTML}{2E6DA4}
\definecolor{codebg}{HTML}{F4F4F4}
\definecolor{codetxt}{HTML}{1A1A1A}

% --- Espacement ---
\setlength{\parindent}{1.5em}
\setlength{\parskip}{4pt}
\singlespacing

% --- Titres sans soulignement ---
\titleformat{\section}{\large\bfseries\color{hdr}}{\thesection}{1em}{}
\titleformat{\subsection}{\normalsize\bfseries\color{hdr}}{\thesubsection}{1em}{}
\titleformat{\subsubsection}{\normalsize\itshape\color{hdr}}{\thesubsubsection}{1em}{}
\titlespacing*{\section}{0pt}{16pt}{6pt}
\titlespacing*{\subsection}{0pt}{10pt}{4pt}
\titlespacing*{\subsubsection}{0pt}{8pt}{3pt}

\newcommand{\tblhdr}[1]{\cellcolor{hdr}\textcolor{white}{\textbf{#1}}}

\begin{document}
\sloppy

\begin{titlepage}

% --- En-tête université ---
\noindent
\begin{minipage}[t]{0.6\textwidth}
  \vspace{0pt}
  {\small\textbf{Université Paris Cité}}\\
  {\small UFR Mathématiques et Informatique}\\
  {\small L3 Informatique --- Intelligence Artificielle}
\end{minipage}
\hfill
\begin{minipage}[t]{0.3\textwidth}
  \vspace{0pt}
  \raggedleft
  \includegraphics[width=\textwidth]{Universite_Paris_logo_horizontal_2000px.png}
\end{minipage}

\vspace{0.5cm}
{\color{hdr}\rule{\textwidth}{1pt}}

\vspace{3cm}

% --- Titre ---
\begin{center}
  {\small\color{gray}\MakeUppercase{Rapport de Projet}}\\[1cm]
  {\fontsize{42}{48}\selectfont\bfseries\color{hdr} La Percée}\\[0.4cm]
  {\large\color{accent}\textit{Implémentation et analyse d'algorithmes de jeu}}\\[1cm]
  {\color{accent}\rule{0.4\textwidth}{0.6pt}}\\[0.8cm]
  {\small\color{gray}
    Minimax \quad$\bullet$\quad Alpha-Beta \quad$\bullet$\quad
    Fonctions d'évaluation \quad$\bullet$\quad Tournoi IA
  }
\end{center}

\vspace{1.5cm}

% --- Date et année au centre ---
\begin{center}
  {\small\color{gray} Année universitaire 2025 -- 2026 \quad|\quad 08 Mai 2026}
\end{center}

\vfill

% --- Filet + auteurs ---
{\color{hdr}\rule{\textwidth}{1pt}}
\vspace{0.4cm}

\noindent
\begin{minipage}[t]{0.5\textwidth}
  \textbf{Auteurs}\\[0.3cm]
  \begin{tabular}{ll}
    KEITA Fodé Laye   & \texttt{22525615} \\
    CRISOSTO Jordan   & \texttt{XXXXXXXX} \\
  \end{tabular}
\end{minipage}

\end{titlepage}

\renewcommand{\contentsname}{Sommaire}
\tableofcontents
\setcounter{section}{0}
\renewcommand{\thesection}{\arabic{section}}

\section{Présentation générale du jeu}

\subsection{Description du jeu}

La Percée (en anglais \textit{Breakthrough}) est un jeu de stratégie abstrait à deux joueurs,
joué sur un plateau de $8 \times 8$ cases. Il a été inventé par Dan Troyka en 2000 et a remporté
le concours de jeux abstraits de l'année. Chaque joueur dispose de 16 pions placés sur les deux
premières rangées de son côté du plateau. Les Blancs occupent les rangées 7 et 8, les Noirs les
rangées 1 et 2.

Ce jeu appartient à la catégorie des jeux à \textbf{information parfaite} et \textbf{à somme nulle} :
les deux joueurs voient l'intégralité du plateau à tout moment, et le gain de l'un correspond
exactement à la perte de l'autre. Il ne comporte aucun élément de hasard.

\subsection{Règles principales}

Les règles de La Percée sont volontairement simples, ce qui le rend facile à prendre en main
tout en offrant une profondeur stratégique réelle.

\begin{itemize}
  \item Les joueurs jouent à tour de rôle. Les Blancs commencent la partie.
  \item Un pion peut \textbf{avancer d'une case vers l'avant}, uniquement si la case est libre.
  \item Un pion peut \textbf{capturer en diagonale avant} (gauche ou droite) une pièce adverse.
  \item Un pion ne peut ni reculer, ni se déplacer latéralement, ni sauter par-dessus un autre pion.
  \item Il n'existe aucune promotion, prise en passant, ni coup spécial.
\end{itemize}

\subsection{Conditions de victoire}

Un joueur remporte la partie dans l'un des trois cas suivants :

\begin{itemize}
  \item Il \textbf{atteint la dernière rangée adverse} avec au moins un de ses pions.
  \item Il \textbf{capture tous les pions adverses}.
  \item L'\textbf{adversaire n'a plus aucun coup légal} disponible --- il est alors immédiatement déclaré perdant.
\end{itemize}

La partie ne peut pas se terminer par un match nul dans les règles standards.

\subsection{Intérêt algorithmique}

La Percée présente plusieurs propriétés qui en font un excellent candidat pour l'étude des
algorithmes d'intelligence artificielle.

Son \textbf{facteur de branchement} est élevé : en début de partie, chaque joueur peut jouer
jusqu'à 30 coups différents (16 pions, chacun pouvant avancer ou capturer). Ce facteur
diminue au fil de la partie à mesure que les pions sont capturés, mais reste suffisamment
élevé pour rendre l'exploration exhaustive de l'arbre de jeu impossible sans élagage.

La \textbf{profondeur de recherche} nécessaire pour jouer correctement est également significative :
un pion doit traverser 6 rangées pour atteindre le camp adverse, ce qui implique des
séquences de coups longues à anticiper.

Ces caractéristiques rendent La Percée particulièrement adaptée à l'algorithme
\textbf{Minimax avec élagage Alpha-Bêta}, dont l'efficacité dépend précisément de la capacité
à élaguer des branches dans un arbre de grande taille.

\subsection{Hypothèses et simplifications}

Dans le cadre de ce projet, nous avons retenu les règles standards de La Percée sans
modification ni simplification. L'interface du jeu est entièrement en ligne de commande :
le plateau est affiché en mode texte, et les coups sont saisis au format algébrique
(par exemple \texttt{a2 a3} pour avancer le pion de la colonne~a, rangée~2 vers la rangée~3).
Aucune interface graphique n'a été développée, conformément aux consignes du projet.

\end{document}